\chapter{Benchmarks}
\section{CFD-DEM}
\subsection{Alternative Benchmarks}

\only<article>{
The benchmarks in the section where discussion as possible CFD-DEM benchmarks but are currently not being actively considered as official ON-DEM benchmarks. 

}

\begin{frame}
\frametitle{Granular Rayleigh-Taylor Instability Benchmark}
\framesubtitle{Status: On Hold}

    \textbf{Benchmark Objective}
    \begin{itemize}
        \item The primary aim is to recover the results from section 7 of \cite{robinson2014fluid}.
        \item It is treated as a pure benchmark against the Robinson data rather than a validation case, as it is unclear why the system should strictly follow the Rayleigh-Taylor instability.
    \end{itemize}

    \vspace{0.3cm}
    \textbf{Physical Setup \& Methodology}
    \begin{itemize}
        \item The simulation involves a set of granular materials arranged on a grid falling into a fluid.
        \item A perturbation is applied to the initial particle positions.
        \item The growth on this interface is then analysed as a function of the perturbation's wave number.
    \end{itemize}
    
    \vspace{0.3cm}
    \alert{Problem: The accuracy of the original data cannot be verified.} 
\end{frame}
\begin{frame}
\frametitle{Sedimentation of a Constant Porosity Block Benchmark}
\framesubtitle{Status: On Hold}

    \textbf{Benchmark Objective}
    \begin{itemize}
        \item The primary target is to recover the analytical solution presented in equation (36) from section 5 of \cite{robinson2014fluid}. 
        \item This case is strictly classified as a verification benchmark. 
    \end{itemize}

    \vspace{0.3cm}
    \textbf{Physical Setup \& Methodology}
    \begin{itemize}
        \item The simulation models a porous block of granular materials falling within a container. 
        \item It serves as a natural extension of a single spherical particle falling in a fluid under a gravitational field, but importantly includes fluid back flow. 
        \item There is a known theoretical, analytic solution for the terminal velocity of the particle for a given applied drag model. 
        \end{itemize}
    
    \end{frame}
    
    \begin{frame}
\frametitle{Segregation in a Fluid-Saturated Rotating Drum Benchmark}
\framesubtitle{Status: On Hold}

    \textbf{Benchmark Objective}
    \begin{itemize}
        \item The target is to validate the numerical codes against experimental data of segregation profiles. 
        \item Unlike the previous cases, this is strictly classified as a validation benchmark. 
    \end{itemize}

    \vspace{0.3cm}
    \textbf{Physical Setup \& Methodology}
    \begin{itemize}
        \item The setup considers two different sizes of glass beads inside a rotating drum filled with a viscous liquid. 
        \item The goal is to analyse the segregation profile as a function of drum RPM, fluid density, and fluid viscosity. 
         \end{itemize}
    

    
    \vspace{0.3cm}
    \alert{Problem: The original experimental paper was never published and the data is missing, though the experimental setup still exists to recreate it.} 
\end{frame}

\section{Industrial Benchmarks}
\begin{frame}
\frametitle{Industrial Challenge: Powder Dosing (Double Screw Feeder)}
\framesubtitle{Status: Active}

    \textbf{Challenge Objective}
    \begin{itemize}
        \item This is an active challenge aimed at evaluating \lq blind\rq\ DEM simulations against experimental results of the dosing process. 
        \item The overarching goal is to enhance precise material delivery while optimising consistency and reducing errors. 
        \item Key comparison parameters include achieved mixing indices (e.g., coefficient of blending performance), segregation amounts, and the effects of agitator and screw rotational speeds. 
    \end{itemize}

    \vspace{0.3cm}
    \textbf{Industrial Partnership \& Data}
    \begin{itemize}
        \item The industrial partner, Bayer, is providing two materials alongside relevant material characterisation and laboratory-scale data.
        \item The provided lab data includes starting mass, filling level, rotational screw and agitator speeds, and dosing time. 
    \end{itemize}
\end{frame}

\begin{frame}
\frametitle{Powder Dosing Challenge: Setup \& Implementation}
\framesubtitle{Status: Active}

    \textbf{Simulation Setup \& Flexibility}
    \begin{itemize}
        \item Participants must calibrate their chosen contact models using the provided material data. 
        \item Because the process involves fine and cohesive particles, participants have the flexibility to choose appropriate contact laws (e.g., Hertzian/Linear elastic with JKR or DMT for adhesion).
    \end{itemize}

    \textbf{Material Properties \& Availability}
    \begin{itemize}
        \item The measured probability density functions (PDFs) of the materials are provided, alongside literature data for paracetamol (e.g., density of 1200 $\text{kg/m}^3$) for guidance. 
        \item Full descriptions of the experiments, equipment geometry, and datasets are hosted on Zenodo \href{[https://zenodo.org/records/18588605](https://zenodo.org/records/18588605)}{Record 18588605}.
    \end{itemize}
    
    \textbf{Participating Codes \& Leads}
    \begin{itemize}
        \item \textbf{MercuryDPM:} \alert{Holger Götz.}
        \item \alert{Note: Every benchmark requires at least two different DEM codes, so additional collaborators are needed.} 
    \end{itemize}
\end{frame}